\section{近期研究（2026 年 8--11 月）}
\paragraph{GaussianSelector：将用户笔迹提升为高斯超点图的前背景约束}
2026.08~\cite{2608_01492}。GaussianSelector 将已重建场景的高斯按外观、不透明度和空间邻近组织为超点图，用光栅时的透射与覆盖把前背景笔迹落到对该像素有贡献的原语，再在图上给其余超点求选择结果。方法直接操作已有高斯，无需为选择任务训练附加网络。

论文在 LLFF-NVOS 的八个任务和 3D-OVS 五个场景比较，并区分单轮与追加笔迹。结果必须结合用户输入视图数和交互轮数解读，遮挡区域能否选中仍取决于可见笔迹与图中连接关系。

\paragraph{Super-Gaussian：以分层分组支持体数据的虚拟现实交互}
2026.08~\cite{2608_04475}。Super-Gaussian 将体可视化的高斯基元按特征聚成可交互单元。在虚拟现实中，用户先用稀疏种子与图传播选取区域，再作组级修正和原语级细化；自然语言界面用于语义标注与查询。

论文将几何选择对比、头显渲染时间与测地 SLIC 消融分开评估。所用初始聚类由预处理得到，窗口参数控制分组粗细，在低对比或边界不明的区域可能产生不稳定分组；图传播也受交互种子质量影响。

\paragraph{OutLangSplat：融合区域结构与语义以理解无人机户外场景}
2026.08~\cite{2608_04560}。OutLangSplat 将二维视觉语言特征和三维高斯邻域结构在区域级对齐，随后以像素上的高斯渲染贡献及跨视图语义一致性加权聚合，减少遮挡和噪声视图的错误响应。

论文在 Buildings1、Buildings2、Polytech 和 Campus 四个无人机户外场景验证开放词汇语义分割与实例定位。它不重新优化场景重建，几何不完整或高斯位置错误会影响对应和聚合，作者将小对象语义与重建联合优化列为后续方向。

\paragraph{CDSeg：用可渲染高斯载体把二维标签回写到三维}
2026.08~\cite{2608_05482}。CDSeg 用已标定视图渲染产出的像素至原语关联，把外部二维掩码的标签融合到可见原语。模式 I 保留输入点索引，便于回写网格或激光雷达点云；模式 II 将标签留在已有、已优化的高斯场景中，可随视图重新渲染。

论文用 DesktopObjects-360 和 NeRDS-360 评价优化载体上的实例任务，用 ScanNet-v2 验证语义标签回写，并在 KITTI-360 的静态车辆窗口考察可见性限制。从未出现在输入视图的表面不会获得标签，已有高斯的空缺和失真也会传递到结果。

\paragraph{InstanceSplat：从无位姿图像前馈重建带实例身份的高斯场}
2026.08~\cite{2608_07144}。InstanceSplat 由无位姿多视图预测共享高斯表示，在同一场景中编码几何、外观、实例和语言语义。跨视图实例对齐为可渲染的实例特征提供监督，实例结构进一步参与联合重建与语义学习。

论文在 ScanNet 与 ScanNet++ 混合训练，在 50 个留出 ScanNet 场景比较 2、4、8 和 16 个输入视图。评估分辨率为 448×448；论文分别报告新视角渲染与分割指标，并指出 4 视图结果中的平均准确率会受采样目标的类别比例影响。

\paragraph{LEGO：把部件、对象与场景的语义按结构层编码}
2026.08~\cite{2608_10057}。LEGO 将多视图 SAM 分割重排成三维一致的结构层级，再分别学习高斯子特征、绑定 CLIP 语言特征并组织为分层场景图，以减少粗层和细层监督互相干扰。场景图还能供语言模型按层检索实体和上下文。

论文在 NVOS 和 SPIn-NeRF 评价提示分割，在 LERF-OVS 与补充的 3D-OVS 等设置检查开放词汇理解，报告 NVOS 的 mIoU 提升最高达 1.6 个百分点。方法依赖初始几何和多视图掩码；开放词汇结果对应论文使用的语义模型和评估场景。

\paragraph{CausalSplat：以高斯语义场与场景图支撑指代推理}
2026.08~\cite{2608_11150}。CausalSplat 先将二维分割提升为高斯语义场，再把实例的多模态属性及空间连接组织成图。视觉语言模型结合图解析复杂指代表达并定位目标，输出可分割的三维实体。

作者构建 Causal-LERF 和 Causal-ScanNet，分别在渲染像素和三维点上计算 mIoU。补充实验拆分空间、常识、可供性和预测或反事实问题。这里验证的是通过这些描述识别目标的表现，输出质量取决于感知和语言推断阶段所获得的证据。

\paragraph{Seed2GS：一次识别并传播对象种子以提取高斯}
2026.08~\cite{2608_11928}。Seed2GS 在冻结的三维场景中，用文本对一张带相机位姿的参考渲染识别目标，借虚拟视图跟踪扩大覆盖，然后通过固定光栅化为各高斯拟合临时前景分数。它可以不使用原始重建相机，但查询仍需给定一张参考视图和该视图的位姿。

论文在用户选择能看见目标的 LERF-MASK 参考视图条件下报告 92.1\% mIoU 和 9.3 秒，并在 3D-OVS 报告 95.7\% mIoU。前一数字未对任意参考视图平均；识别错误、严重遮挡、缺少几何和长距离跟踪漂移都可能导致失败。

\paragraph{QAGaussian：用关系推理选取自由指代表达指定的高斯}
2026.08~\cite{2608_16103}。QAGaussian 从语言查询产生多尺度候选槽，再通过候选关系图和几何约束组合证据，输出逐高斯的掩码概率。例如带参照物的表达中，空间关系可用于减轻目标与参照物被一起选中的问题。

作者在 Mosaic3D-5.6M 预训练并用该数据的验证集选定设置，随后测试 ScanRefer、ReferIt3D、Multi3DRefer、ReferSplat 和部件数据。不同标注被转换到共同评价设定，比较结果应按这些对应协议解释；预测仍依赖高斯重建和指代数据的覆盖。

\paragraph{GroupForward：以实例分组承载前馈场景语言语义}
2026.08~\cite{2608_17535}。GroupForward 从稀疏、无位姿且未标定图像前馈重建高斯，以紧凑实例嵌入建立跨视图分组。开放词汇特征按实例聚合，随后以实例构建场景图并结合多视图观察，供视觉语言模型搜索目标和输出指代分割。

作者以预训练 VGGT 初始化几何部分，并使用 8 维实例嵌入。论文将语义三维重建与场景指代推理分开评价；低维只描述实例分组嵌入，语言特征仍在实例层维护。图推理能否定位目标还取决于前馈阶段能否把目标分成有效实例。

\paragraph{LetOccVote：利用跨帧投票改进弱监督三维占据预测}
2026.09~\cite{2609_04846}。LetOccVote 的深度投票把多帧在共同空间里的支持用于提升和深度监督，再通过高斯将多帧信息传给三维占据表示。语义投票用同意和冲突区域筛选不可靠的二维伪语义监督。

在 Occ3D-nuScenes 的主表，较大模型报告 53.27 IoU、20.39 mIoU，另按论文标星类别集合报告 23.11 mIoU。两个语义数值使用不同的类别口径。方案借助跨观察共识减轻伪标签噪声，评价以多帧驾驶场景和二维伪标签监督条件为背景。

\paragraph{LangStreet：把语言证据绑定到锚点解码的持久地址}
2026.09~\cite{2609_11616}。LangStreet 将随视图生成的高斯原语用作路由，把多视图语义观察累加到持久槽位并向上汇成锚点摘要。对于观察稀疏的槽位，通过保持锚点方向的层级补全增强稳定性，再用相对锚点编码保存局部差异。

作者在 KITTI-360、Virtual KITTI 2 和 Waymo 上共享冻结的场景重建、解码器及二维语言引擎进行比较。论文显示紧凑字段可以接近完整字段并节省存储；结论对应同一预处理观察和持久结构，不能解释为重新训练重建或二维模型后的端到端提升。

\paragraph{MoQSplat：将高斯场景按空间和质量层送入渐进传输}
2026.09~\cite{2609_18624}。MoQSplat 在 Media over QUIC 的内容组织中，将场景空间区域映射到 Track，再以语义分组和由粗到细的质量层组织对象。接收端按关注区域及层级决定请求优先级，发布者和转发者无需为各接收方维护单独的选层状态。

论文先测内容准备质量，再测网络条件变化下的端到端接收。两组使用的共同参照是 Mip-NeRF 360 的 Bicycle，内容准备以固定内参、距离和高度的绕行轨迹计算 VMAF。因此不透明度排序和带宽适配的观察属于单场景初步实验。

\paragraph{NormLift：用跨视图支持校准高斯语言特征置信度}
2026.09~\cite{2609_18898}。NormLift 将单位球面上的逐高斯余弦对齐写成可闭式求解的提升问题，结果对应反投影特征的归一化。聚合前的向量范数反映观察的一致程度，再以有效视图支持校准，作为邻域模式投票修正低可信原语的依据。

作者以 OpenGaussian 协议在 ScanNet 测试开放词汇三维分割并比较 mIoU 和 mAcc。后处理以既有高斯场景为基础，置信度与投票利用观察和邻域的证据，不补回原场景缺失的表面。

\paragraph{ParticleSplat：由多视图自监督学习可编辑的对象级潜粒子}
2026.09~\cite{2609_19463}。ParticleSplat 从已知相机姿态和内参的观察学习结构化三维潜粒子，并将粒子前馈解码为显式高斯。粒子形成对象级控制接口，改变对象潜变量或位置后可观察对应的新视图变化。

论文把分解、新视图渲染和操作学习分为不同实验，其中真实桌面场景使用手持手机捕捉多视图，机器人决策收益另有仿真设置。真实捕捉无需固定多机阵列仍需获得带位姿的多视图，不能理解成相机位姿无关。

\paragraph{CoRef-GS：通过地图配准和实例融合做跨机器人语言指代}
2026.09~\cite{2609_20586}。CoRef-GS 先构建可按实例查询的局部语义高斯地图，通过跨机器人配准和实例融合将观察统一到共享坐标，再按查询视角使用空间关系选出语言指定的对象。

作者在含真实和模拟室内场景的 CoQuad-Ref 上区分地图配准、开放词汇查询和协作指代结果。主结果报告粗 Sim(3) 初始化已改善配准，完整方案在模拟设置中进一步降低误差。证据来自互补观察的室内四足机器人设置，不代表大规模动态车队的通信和时间同步表现。

\paragraph{建筑感知卫星高斯重建：经掩码质检和区域误差选择权重}
2026.09~\cite{2609_25578}。该工作用建筑掩码调节卫星高斯重建的区域损失，控制器同时参考掩码质检、场景元数据和验证误差决定是否采用建筑先验。系统记录评价结果，使数字表面模型的选择可追溯到场景条件。

论文用注册后的 DFC2019 风格 DSM 与语义标注计算建筑、植被和全场 MAE。JAX\_004 上的建筑权重 2.0 将建筑区域误差从 0.844 米降至 0.806 米，但分阶段方案在全场误差上的折中更好。改善以通过掩码质检的场景为条件，不能只用单一建筑误差决定部署。

\paragraph{PePESeg3D：在几何和特征两阶段注入感知先验}
2026.09~\cite{2609_28645}。PePESeg3D 首先以单目深度和分割掩码引导高斯几何贴近结构边界，再通过深度颜色线索与跨视图质心约束学习多尺度语义，以补充二维掩码在部分粒度上的缺漏。

论文在 SPIn-NeRF、LERF-Mask 等设置报告标签传播与提示分割结果，另用 PSNR、SSIM 和 LPIPS 评价重建，两个任务的指标不能互代。作者观察到共享特征空间的尺度门控会使不同粒度的掩码互相影响，尤其在邻近尺度检索时。

\paragraph{SplatLabel：以含生命周期的动态高斯生成三维语义伪标签}
2026.09~\cite{2609_29836}。SplatLabel 将连续运动、生命周期和存在边界纳入动态高斯属性，让二维语义线索与激光雷达几何在时间上共同决定三维语义和体积占用。高斯在此是时空伪标签生成器。

论文以 SemanticKITTI 验证序列 08 评价，使用 SAM 3 和类别名称直接生成语义特征。方法依赖精确相机和 LiDAR 同步且显存开销较大，恶劣光照、天气或传感器质量退化会影响结果。

\paragraph{EviSplat：把多视图语义证据保留到查询时再聚合}
2026.09~\cite{2609_34853}。EviSplat 保留实例的各视图特征，并为高斯学习反映不同可见外观的紧凑证据场。查询时先按相关性选取实例观察，再用原语局部的可靠证据细化分数，不需要为了新查询重训表示。

论文在 LERF-OVS 和 3D-OVS 评价渲染掩码，在 ScanNet 按 OpenGaussian 设定评价 10 个场景的三维点，并将渲染后查语言特征与在三维先选高斯的方案分列。结论需区分这两类输出及其协议；保留逐视图证据也改变了语义字段的存储结构。

\paragraph{Imagine3D-LLM：把紧凑高斯重建变成多模态问答的训练任务}
2026.09~\cite{2609_38177}。Imagine3D-LLM 在多视图视觉标记后加入可学习的高斯摘要标记，把它们解码为紧凑场景表示，并与语言预测共同训练。对多图问答而言，高斯在这里提供显式空间结构监督，促进跨视图统一物体表征。

论文比较多类空间推理与三维理解基准，还分析生成高斯和内部特征。作者指出不用紧凑高斯教师也能学习摘要标记，但达到相近下游结果需要更多训练步骤。问答收益与场景几何的细节准确度需要分别度量。

\paragraph{MEGA：从分割高斯的虚拟视图蒸馏对象级网格}
2026.10~\cite{2610_01707}。MEGA 从已分割的高斯对象渲染一组虚拟视图，对每个对象以光度和掩码监督蒸馏神经表面，输出可用于对象级操作的网格。分割和表面蒸馏承担不同环节，蒸馏不会绕过初始对象分离。

作者的定量设置使用三排、每排 72 个虚拟视点。在 LERF-Mask 掩码的额外实验中，MEGA 平均 mIoU 为 90.1，作为分割教师的 SAGO 为 91.0，不能把该指标描述成超过教师。网格提取质量与该分割和虚拟视图采样有关。

\paragraph{MaRO-GS：以跨视图掩码可靠性约束对象重建}
2026.10~\cite{2610_06472}。MaRO-GS 根据几何一致性过滤不可靠对象掩码，再限定高斯密度控制只支持具有对象证据的位置，以抑制背景高斯混入，轮廓训练为对象边界提供约束。

主渲染比较采用 LERF-Mask、Mip-NeRF 360 和 Tanks and Temples 上统一的掩码输入协议，测量与掩码目标的渲染一致性。作者指出严重错误的掩码仍会破坏结果；Hard DTU 子集较小，保守过滤的召回只有 4.6\%，不能将整体改善全归因于过滤。

\paragraph{OpenSplatGraph：从可靠性语义高斯场按需生长场景图}
2026.10~\cite{2610_07569}。OpenSplatGraph 为稠密语义高斯地图保存可靠性信息，按语言查询提取实例并将其关联到现有节点。持久的场景图逐步维护对象身份、属性和关系，为随后查询保留可复用的结构。

实物实验使用搭载 RGB-D 相机的 Stretch 3 移动机械臂，在探索中融合语义地图，另有三维分割和场景图评价。该框架以查询触发实例提取，不维护独立于查询的完整对象清单；拥挤和重遮挡仍会使轻量关联变得不可靠。

\paragraph{训练后语义提升：分别估计检测器、提升和表示误差}
2026.10~\cite{2610_08756}。该方法不重训既有高斯。给定单个目标类别，它累积各视图的目标和非目标证据，以高阈值选择种子，再以较低阈值接受含种子的连通区域。评估将预测映射到相机可见且有标注的网格顶点，从而排除不可见表面带来的混淆。

论文分别比较标注掩码和检测器掩码，在 Replica 选定设置并在 ScanNet++ 验证迁移。类若未出现在任何输入掩码就没有可提升的证据，且逐类处理可能让同一高斯落入多个类别，不能将该方法直接视为单次多类别分割。

\paragraph{可测量语义高斯：面向建筑巡检的区域和中心深度监督}
2026.10~\cite{2610_13143}。该工作在高斯优化中联合使用目标区域约束与逐点深度监督，直接约束三维高斯中心的位置，并构建语义场，使建筑高斯既可观察部件类别，也能从显式结构中取点测量。

作者分别在合成 QuakeCity 和手机采集的 HKUPB 上评价区域 RGB、中心几何与语义表现。中心深度监督改善的是总体对齐，局部仍可能因原语稀疏影响精确取点。作者指出细小裂缝和局部结构很难重建，同时输入位姿和深度仍依赖摄影测量。

\paragraph{DIGS：在高斯上保留开放词汇标签分布}
2026.10~\cite{2610_16177}。DIGS 用逐高斯的前若干候选标签分布替代单一语义向量，并以减温 Dirichlet 更新限制单帧观察的影响，减轻来自多实体边界的语义混合。实例级视觉和文本表征帮助在不同粒度下给高斯赋分。

论文以四个开放词汇基准测试，其中 LERF-OVS 的四个场景按 OpenGaussian 三维原生协议评价实例定位与分割。映射时间、显存和单次查询另行报告，方法依赖由基础分割和语言模型得到的输入证据。

\paragraph{目标感知高斯语言地图：最小特征注入与中心化定位}
2026.10~\cite{2610_21619}。该方法将指定目标在各视图中的视觉特征向文本嵌入靠拢，以闭式求出的最小强度注入，再将更新后的特征提升至高斯地图。查询时的中心化读出消除注入产生的共模分量，用以返回目标三维位置。

论文在三个基准将开放词汇分割与语言目标定位分开评价，并展示真实四足机器人。未做中心化的消融在三个数据上均可能因查询分数高度相关而退化；例如 ScanNet 的首选定位从基线 72.9 降至 30.0。论文的导航位置比较固定底层地图，不能等同于不同系统完整导航管线的比较。

\paragraph{O-SCD++：以递归统计在固定高斯地图检测变化}
2026.10~\cite{2610_24462}。O-SCD++ 利用高斯线性合成，将由新观察获得的变化证据累积到固定参考场景的逐高斯状态中。递推更新和闭式决策替代反复缓存、重放观测并迭代优化的处理方式。

作者以真实多视图场景变化的 PASLCD 为主要实验，场景包含多个几何和表面变化，并有不同光照的重访数据。系统无法恢复输入线索未识别到的变化；例如只用语义线索时，同类表面的颜色改变可能被忽略。

\paragraph{PanoGraphGS：统一高斯全景分割和楼层至对象的层级场景图}
2026.10~\cite{2610_27696}。PanoGraphGS 在共享的高斯几何上训练类别原型和持久实例关联，从同一表示解码新视角全景掩码与原语级语义。语义及空间结构进一步组织成楼层、房间、背景实体与对象的层次图。

作者使用 ScanNet++、Replica 和自采 NaviGS 验证，论文将二维全景分割与三维场景图的楼层、房间和对象节点准确率分开报告。开放式厨房一类缺少显式边界的空间难以划分房间；图构建在当前版本中离线运行。

\paragraph{自动编辑参考高斯：为在线贝叶斯变化检测合成证据}
2026.10~\cite{2610_28490}。该工作用自动编辑后的参考高斯生成可监督视图，预先拟合场景专用的变化证据模型。上线后以局部贝叶斯因子和递推的历史证据估计变化，无需对在线数据反向传播或重放观测。

实验把新占据与腾空区域都纳入变化掩码，以真实 PASLCD 和合成 CCS3D 为比较环境。作者报告 PASLCD 的前景 mIoU 达 67.64\% 并以 17.2 帧每秒处理。简单对象编辑与重光照覆盖不了所有真实变化；证据累积在固定参考几何上，限制了新出现对象的位置表达。

\paragraph{R4DGS：为动态高斯场构建可持续指代的实体库}
2026.10~\cite{2610_30326}。该工作将动态场景中的语言指代表述为目标实体选择。ReferGaussian 从预计算的四维高斯及多帧视觉掩码构建持久实体，存入 EntityBank，以多模态模型解释短语与时序约束，并在实体层输出结果，避免另训逐场景语义场。

评价使用 R4D-Bench-QA 和四个场景的 4D LangSplat 协议，把时序状态、多目标与无目标指代纳入测试，并对比语言场和分割先行路线。方法依赖预计算的动态场景与掩码轨迹，无法单靠实体推理恢复底层场景未观察或未重建的内容。

\paragraph{RecallGS：用当前几何核实可回用的历史语义观察}
2026.10~\cite{2610_31293}。RecallGS 从历史图像检索与当前查询相似的帧，再通过透明度合成中的像素贡献检查历史帧是否真正支持当前高斯，并只融合满足几何支持的线索。它将记忆用于语义特征补充，而不需再训练新的语义网络。

论文沿用 EmbodiedSplat 的预处理测试场景，在 ScanNet 分 10、15、19 类评估，ScanNet++ 采用 20 类设置。结果限于这些室内地图和已有历史视图；检索得到视觉相似帧不等于看到了同一表面，几何验证的用途正是过滤此类误配。

\paragraph{GEMS：保存高斯变化量并按时间重渲染场景记忆}
2026.10~\cite{2610_33158}。GEMS 针对反复访问的半静态环境，先建高斯地图，再只记录检测到变化的区域，以节省各时点的场景存储。问答代理按问题选择历史时间和相机位姿渲染场景，逐步收集支持答案的视觉证据。

GEMS-Bench 由 3 个真实室内场景与 12 个模拟 HSSD 场景构成，分别含 244 和 229 对问答。问答用模型裁判进行分级。实验基于离散重访和半静态布置，不检验连续快速运动的重建或时间跟踪。

\paragraph{RelateGS：筛选高斯语义观察并回查关系证据}
2026.10~\cite{2610_33868}。RelateGS 在既有高斯重建和带位姿 RGB-D 观察上工作。它在提升语义时过滤弱证据，并以持久对象图汇总多视图信息。当支撑某一关系的观察不足时，系统从既有视图补查相关线索，减少因单视图误判造成的指代错误。

指代实验使用 ReferIt3D 的 Sr3D+ 和 Nr3D，在沿用先前协议的 8 个 ScanNet 室内场景上比较。该机制依赖能正确关联对象和获取支持关系的帧；回用已有图像不能补足从未被观察的线索。

\paragraph{LegoOcc++：共享前馈高斯的占据和开放词汇语义推理}
2026.10~\cite{2610_34093}。LegoOcc++ 从输入图像前馈预测带语言特征的高斯。基元位置、协方差和不透明度经 Poisson 式聚合估计体素占据，并由原语嵌入预测占据体素的类别。实例一致性与共视支持用于减少同一物体内部的语义破碎。

论文分别讨论室内单目和室外驾驶设置。室外方案可以不使用占据标注，室内实验仍需要二值占据监督；不能将二者统称为无需三维几何监督。作者指出视觉与几何差异大的概括类别如家具仍较难统一校准。

\paragraph{ClusterSplat：将高斯指代转为场景适应的离散簇选择}
2026.10~\cite{2610_36462}。ClusterSplat 学习带实例信息的高斯特征，以描述长度准则生成场景适应的簇，再按语言查询为各簇排序。最终可渲染被选簇的分数产生二维掩码，从而把按单个原语过滤改为以候选对象为单位选择。

论文在配对 ScanNet v2 与 ScanRefer 的场景以及 Ref-LERF 的二维协议比较。作者指出主要瓶颈是候选簇：相邻同类粘连或把一个对象拆成数簇都可能让首选验证通过但空间 IoU 仍不足。

\paragraph{按查询分割高斯：从一个二维枢轴传播并投票到三维}
2026.11~\cite{2611_19429}。该方法以密集视觉语言特征给文本查询生成相似度图，选择二维枢轴区域，随后以视频分割器在有限视图预算内跟踪并细化掩码。根据每个高斯的可见证据投票，形成此次查询对应的三维实例掩码。

论文在 LERF-OVS 将首次查询所需时间、分割 mIoU 与定位准确率分别报告，并另列多个室内基准。该方案使用预训练模型和图像特征处理；视点变化大时跟踪可能漂移，严重遮挡也会使最终掩码不完整。

\paragraph{VODA：从屏幕掩码监督扩展到深度簇内的全部高斯}
2026.11~\cite{2611_21383}。VODA 将同一像素参与渲染的高斯按局部深度分组，在目标深度簇中直接监督每个高斯的特征，避免训练只惩罚对渲染掩码贡献较大的表面原语。该目标可以作为现有语言和实例分割方案的特征学习组件。

论文以 Removal-3D 的移除前后深度变化度量检验体积选择完整性，并在 SPIn-NeRF 和 360-USID 的成对实拍中验证遮罩质量。实拍深度差参考由独立拟合的有对象与无对象高斯场构造，作为代理指标使用，不能视为真值三维几何。
